\documentclass[11pt]{article}
\usepackage[margin=0.75in]{geometry}
\usepackage[T1]{fontenc}
\usepackage[utf8]{inputenc}
\usepackage{lmodern}
\usepackage{microtype}
\usepackage{enumitem}
\usepackage{booktabs}
\usepackage{longtable}
\usepackage{array}
\usepackage{tabularx}
\usepackage{xcolor}
\usepackage{hyperref}
\usepackage{listings}
\usepackage{parskip}
\usepackage{multicol}
\definecolor{linkblue}{RGB}{0,0,140}
\definecolor{revisionnote}{RGB}{130,80,0}
\definecolor{diffgreen}{RGB}{220,247,226}
\definecolor{diffgreentext}{RGB}{20,92,45}
\definecolor{diffred}{RGB}{252,226,226}
\definecolor{diffredtext}{RGB}{145,30,30}
\definecolor{diffwhite}{RGB}{255,255,255}
\definecolor{diffunchangedtext}{RGB}{55,55,55}

\newcommand{\revshort}[1]{\texttt{#1}\textsuperscript{\textcolor{revisionnote}{\dag}}}
\newcommand{\revmark}{\textsuperscript{\textcolor{revisionnote}{\dag}}}

\newcommand{\diffadd}[1]{%
  \begingroup
  \setlength{\fboxsep}{2pt}%
  \colorbox{diffgreen}{\textcolor{diffgreentext}{\strut #1}}%
  \endgroup
}
\newcommand{\diffremove}[1]{%
  \begingroup
  \setlength{\fboxsep}{2pt}%
  \colorbox{diffred}{\textcolor{diffredtext}{\strut #1}}%
  \endgroup
}
\newcommand{\diffsame}[1]{%
  \begingroup
  \setlength{\fboxsep}{2pt}%
  \colorbox{diffwhite}{\textcolor{diffunchangedtext}{\strut #1}}%
  \endgroup
}
\hypersetup{colorlinks=true,urlcolor=linkblue,linkcolor=linkblue,pdftitle={KeyQuorum File History Architecture},pdfauthor={Bailey Forbes}}
\lstset{
  basicstyle=\ttfamily\small,
  frame=single,
  breaklines=true,
  breakatwhitespace=false,
  columns=fullflexible,
  keepspaces=true,
  showstringspaces=false,
  escapeinside={(*@}{@*)}
}
\title{KeyQuorum File History Architecture\\\large Topology-Aware Provenance, Revision Trust, and Comparison to Git}
\author{Bailey Forbes}
\date{September 2026}
\begin{document}
\maketitle
\tableofcontents
\newpage

\section*{\textcolor{revisionnote}{\dag\ Revision Notation Note}}
\addcontentsline{toc}{section}{Revision Notation Note}

Throughout this document, compact labels such as \revshort{R1}, \revshort{R2}, \revshort{R3A}, and \revshort{R3B} are used only as explanatory shorthand in diagrams and examples.

They are \textbf{not} production revision identifiers and they are \textbf{not} the literal generated revision label. Each real KeyQuorum revision has three separate representations:

\begin{enumerate}[leftmargin=1.8em]
  \item \textbf{Authoritative revision ID}: a cryptographic SHA-256 identifier derived from the canonical revision object and its parent revision ID(s).
  \item \textbf{Always-generated revision label}: \texttt{R<normalized-file-name>-<UTC timestamp>-<HCP label>}.
  \item \textbf{Optional user label}: a human description analogous to a Git commit message.
\end{enumerate}

For example, the shorthand \revshort{R2} might refer to a revision whose actual values are:

\begin{lstlisting}
revision_id:
    bc2810f4a89d7...

generated_label:
    Rsales-report-20261002T143205.482Z-M.A

user_label:
    Updated sales totals after accounting review
\end{lstlisting}

The dagger symbol (\textcolor{revisionnote}{\dag}) marks shorthand in both prose and code-style examples. In \texttt{lstlisting} blocks it is inserted with the listing escape sequence, so labels such as \texttt{R1}, \texttt{R2}, and \texttt{R3A} visibly carry the same notation marker.

Whenever cryptographic identity, verification, transport, signing, parent linkage, or policy evaluation matters, the full \texttt{revision\_id} is authoritative.

\section{Purpose}
The KeyQuorum File History module is a tamper-evident, topology-aware provenance system for tracked files. It records not only that an action occurred, but also who performed it, which revision was affected, what cryptographic proof was present, how the actor related to the file's Hierarchical Cryptographic Principal (HCP) scope, and what KeyQuorum policy decided.

It should answer questions such as who edited or signed a revision; whether a supervisor countersignature was required; whether a bridge or quorum was used; whether a file was shared with or without a trusted content signature; which revision was actually delivered; whether a newer revision was rejected; who attempted access after expiry; whether password/PIN gates succeeded; and whether history has been modified, replayed, or forked.

The authoritative history travels with the tracked file. SQLite may cache and index the history, but it is not the provenance authority.

\section{Core Architectural Principle}
A tracked KeyQuorum file is:
\begin{center}
\textbf{Stable File Identity + Content + Revision Lineage + Cryptographic Proofs + HCP Topology Context + Policy Decisions + Append-Only History}
\end{center}

\begin{lstlisting}
+----------------------------------------------+
|              KeyQuorum Tracked File          |
+----------------------------------------------+
| Stable File Identity                         |
| File Policy / HCP Scope                      |
| Current Revision                             |
| Trusted Revision Checkpoints                 |
| Cryptographic Proofs                         |
| Embedded History Chain                       |
| History Integrity Root                       |
+----------------------------------------------+
\end{lstlisting}

\section{KQTF: KeyQuorum Tracked File}
Introduce \texttt{KQTF} as the tracked-file container. A logical \texttt{sales-report.xlsx} may internally be represented as \texttt{sales-report.kqtf}, while the UI continues to display the original filename.

KeyQuorum should not append proprietary bytes directly to arbitrary formats such as PDF, DOCX, XLSX, PNG, JPG, or ZIP. Instead KQTF wraps the native payload:
\begin{lstlisting}
KQTF
|-- metadata
|-- native payload
|-- revision checkpoints
|-- signatures / proofs
`-- history
\end{lstlisting}

Editing becomes:
\begin{lstlisting}
sales-report.kqtf
        | checkout
        v
sales-report.xlsx
        | external edit
        v
keyquorum file checkin
        |
        v
sales-report.kqtf
\end{lstlisting}

\section{Tracking Activation}
The default activation point is the first trusted KeyQuorum content signature:
\begin{lstlisting}
Native File
    | first content signature
    v
Generate stable file_id
    v
Create Revision R1(*@\revmark@*)
    v
Sign R1(*@\revmark@*)
    v
Create KQTF
    +--> TRACKING_STARTED
    `--> REVISION_SIGNED
\end{lstlisting}

\section{Stable File Identity}
A tracked file ID is independent of filename, path, device, current holder, current revision, or organizational placement.
\begin{lstlisting}
file_id:       128-bit random identifier
logical_name:  sales-report.xlsx
\end{lstlisting}
Renaming or moving the file does not create a new identity.

\section{Revision Identity, Parent Pointers, and Labels}
Each checked-in content state becomes an immutable revision object identified by a cryptographic \texttt{revision\_id}. The revision ID is the authoritative identifier. It is derived from the canonical revision object and therefore commits to the revision's parent revision ID or IDs, content commitment, author identity, HCP label, UTC creation time, topology generation, policy context, mandatory generated label, and optional user label.

For a normal edit, the revision contains one parent revision ID. A merge contains two or more parent revision IDs. The first revision contains no parents.

\begin{lstlisting}
revision_id =
    SHA-256(
        "KQ-FILE-REVISION-v1"
        || file_id
        || parent_revision_ids
        || content_commitment
        || generated_label
        || optional_user_label
        || author_identity
        || author_hcp_label
        || created_at_utc
        || topology_generation
        || policy_hash
    )
\end{lstlisting}

The parent pointers create a Git-like immutable revision DAG:

\begin{lstlisting}
a41c9e7...  Initial revision
     |
     v
bc2810f...  Accounting update
     |
     v
f93ad44...  Approved correction
\end{lstlisting}

Every revision always receives a deterministic human-readable generated label:

\begin{lstlisting}
R<normalized-file-name>-<UTC timestamp>-<HCP label>
\end{lstlisting}

For example:

\begin{lstlisting}
Rsales-report-20261002T143205.482Z-M.A
\end{lstlisting}

The normalized filename is derived from the logical filename for readability: the extension is omitted, whitespace is converted to hyphens, unsupported punctuation is removed or normalized, repeated separators are collapsed, and the HCP label retains its dotted form. The timestamp is always UTC and should include sub-second precision sufficient to make adjacent operations distinguishable.

A revision may additionally carry an optional user-defined label analogous to a Git commit message:

\begin{lstlisting}
generated_label:
    Rsales-report-20261002T143205.482Z-M.A

user_label:
    Updated sales totals after accounting review
\end{lstlisting}

The generated label is never replaced by the optional user label. Both may be displayed. Because both are part of the immutable revision object, changing the optional label after check-in creates a different revision ID rather than mutating an existing revision.

The generated label and optional user label are descriptive metadata, not authoritative identifiers. Security-sensitive references always use the full cryptographic \texttt{revision\_id}; user interfaces may display an abbreviated hash.

KeyQuorum must also distinguish current from trusted revision:

\begin{lstlisting}
bc2810f...  SIGNED, TRUSTED
f93ad44...  EDITED, UNSIGNED, UNTRUSTED

current_revision_id = f93ad44...
trusted_revision_id = bc2810f...
\end{lstlisting}

A newer byte sequence is not automatically a newer trusted version. In prose, shorthand should be written as \revshort{R1}, \revshort{R2}, or \revshort{R2A}. Listings and diagrams use the same dagger marker through the configured \texttt{lstlisting} escape sequence. These aliases refer to the complete revision object, not merely to the generated label, and are never production revision identifiers.

\section{HCP Topology Integration}
A tracked file may bind itself to an HCP scope, for example:
\begin{lstlisting}
M
`-- M.A
    |-- M.A.1
    `-- M.A.2

file_scope_root = M.A
\end{lstlisting}
This makes \texttt{M.A} the scope owner, \texttt{M.A.1}/\texttt{M.A.2} descendants, \texttt{M} an ancestor, and \texttt{M.S}/\texttt{M.S.1} cross-branch principals.

The history module should reuse shared ancestry primitives rather than separately parse dotted labels:
\begin{lstlisting}
pub fn is_ancestor(ancestor: &str, descendant: &str) -> bool;
pub fn is_descendant(descendant: &str, ancestor: &str) -> bool;
pub fn ancestry_distance(ancestor: &str, descendant: &str) -> Option<usize>;
pub fn lowest_common_ancestor(left: &str, right: &str) -> Option<String>;
\end{lstlisting}

\section{Revision Authority Classification}
\begin{lstlisting}
pub enum RevisionAuthority {
    ScopeOwner,
    Descendant { ancestor: String, depth: usize },
    Ancestor { depth: usize },
    CrossBranch { common_ancestor: Option<String> },
    Unrelated,
}
\end{lstlisting}
This relationship is an input to policy, not merely metadata.

\section{Supervisor Countersignature Integration}
The existing supervisor/delegated-authority model should be reused. If \texttt{M.A.1} edits a document scoped to \texttt{M.A}, policy may require author signature plus direct-parent countersignature:
\begin{lstlisting}
M.A.1 creates R2(*@\revmark@*)
       v
M.A.1 signs R2(*@\revmark@*)
       v
Resolve direct parent: M.A
       v
M.A countersigns
       v
R2(*@\revmark@*) becomes TRUSTED
\end{lstlisting}

History:
\begin{lstlisting}
EDIT_CHECKED_IN
actor: M.A.1
revision: R2(*@\revmark@*)

REVISION_SIGNED
actor: M.A.1
revision: R2(*@\revmark@*)

COUNTERSIGNATURE_ADDED
actor: M.A
for_actor: M.A.1
revision: R2(*@\revmark@*)
relationship: DIRECT_PARENT

POLICY_DECISION
revision: R2(*@\revmark@*)
result: TRUSTED
reason: AUTHOR_SIGNATURE + REQUIRED_PARENT_COUNTERSIGNATURE
\end{lstlisting}

A scope owner such as \texttt{M.A} may self-sign when policy permits.

\section{Configurable Revision Policy}
\begin{lstlisting}
scope_root = M.A

edit_policy:
    descendants_allowed = true

signature_policy:
    scope_owner = SELF_SIGN
    descendants = AUTHOR_SIGN + DIRECT_PARENT_COUNTERSIGN
    ancestors = SELF_SIGN
    cross_branch = AUTHOR_SIGN + BRIDGE_OR_SCOPE_OWNER_APPROVAL
\end{lstlisting}
Policies should be versioned or hashed so later verification knows what rules applied.
The cross-branch rule may instead require both approvals. KeyQuorum exposes
this stricter conjunction as \texttt{author+bridge+owner}: the author's content
signature, a revision-specific bridge approval, and the scope owner's
countersignature must all verify; neither approval substitutes for the other.

\section{Cross-Branch Collaboration and Bridges}
A revision from outside the file scope is explicitly cross-branch. A policy may require a revision-specific bridge approval and scope-owner approval:
\begin{lstlisting}
M.S.1 signs proposed revision
        v
Authorized bridge signs approval of this revision
        v
M.A countersigns
        v
Revision becomes trusted
\end{lstlisting}

Private bridge details must not leak through transferable history. Depending on privacy policy, either omit the bridge event or record only a non-identifying proof outcome.

\section{Signature Classes}
KeyQuorum must distinguish content trust from transport authentication. Recommended classes are:
\begin{itemize}[leftmargin=1.5em]
  \item \texttt{CONTENT\_SIGNATURE}
  \item \texttt{TRANSPORT\_SIGNATURE}
  \item \texttt{COUNTERSIGNATURE}
  \item \texttt{BRIDGE\_SIGNATURE}
  \item \texttt{HISTORY\_EVENT\_SIGNATURE}
\end{itemize}

A content signature may bind:
\begin{lstlisting}
KQ-FILE-REVISION-v1
|| file_id
|| revision_id
|| content_commitment
|| signer_identity
|| topology_generation
|| policy_hash
\end{lstlisting}

A supervisor countersignature may bind:
\begin{lstlisting}
KQ-FILE-COUNTERSIGN-v1
|| file_id
|| revision_id
|| author
|| author_signature_hash
|| supervisor
|| topology_generation
|| policy_hash
\end{lstlisting}

\section{Topology Generation Binding}
Authority is evaluated against the topology that existed at authorization time. If \texttt{M.A.1} later moves departments, an earlier valid signature does not silently change meaning.
\begin{lstlisting}
REVISION_SIGNED
actor: M.A.1
revision: R2(*@\revmark@*)
topology_generation: 7

COUNTERSIGNATURE_ADDED
actor: M.A
relationship: DIRECT_PARENT
revision: R2(*@\revmark@*)
topology_generation: 7
\end{lstlisting}

\section{Trusted Revision Evaluation}
Conceptually:
\begin{lstlisting}
TrustedRevision =
    ValidHistoryChain
    AND ValidContentSignature
    AND ValidTopologyRelationship
    AND RequiredSupervisorApproval
    AND RequiredBridgeApproval
    AND FilePolicySatisfied
\end{lstlisting}
Not every policy requires every component.

\section{Revision Selection for Sharing}
When the newest working revision is untrusted, KeyQuorum should deliver the newest eligible trusted revision or deny delivery.
\begin{lstlisting}
pub struct DeliveryDecision {
    pub candidate_revision: RevisionId,
    pub delivered_revision: RevisionId,
    pub decision: DeliveryDecisionKind,
    pub proof: Option<ContentProof>,
}

pub enum DeliveryDecisionKind {
    CurrentTrustedRevision,
    LastTrustedRevision,
    RequesterAlreadyCurrent,
    DeniedNoTrustedRevision,
    DeniedPolicy,
}
\end{lstlisting}

\section{Sales Report Example}
On October 1, M signs \revshort{R1} and shares it with M.A:
\begin{lstlisting}
TRACKING_STARTED  actor: M    revision: R1(*@\revmark@*)
REVISION_SIGNED   actor: M    revision: R1(*@\revmark@*)
POLICY_DECISION   revision: R1(*@\revmark@*) result: TRUSTED
SHARE_DELIVERED   from: M     to: M.A revision: R1(*@\revmark@*)
\end{lstlisting}

On October 2, M.A edits the file but does not sign it:
\begin{lstlisting}
EDIT_CHECKED_IN
actor: M.A
base_revision: R1(*@\revmark@*)
new_revision: R2(*@\revmark@*)

R2(*@\revmark@*) = CURRENT
R1(*@\revmark@*) = LAST TRUSTED
\end{lstlisting}

When M requests the latest version:
\begin{lstlisting}
FILE_REQUESTED
requester: M
holder: M.A
requested_revision: latest

SHARE_ATTEMPTED
from: M.A
to: M
candidate_revision: R2(*@\revmark@*)
content_signature_present: false

POLICY_DECISION
result: FALLBACK_TO_LAST_TRUSTED_REVISION
candidate_revision: R2(*@\revmark@*)
delivered_revision: R1(*@\revmark@*)
reason: CONTENT_SIGNATURE_REQUIRED

SHARE_DELIVERED
from: M.A
to: M
revision: R1(*@\revmark@*)
\end{lstlisting}
If M.A signs \revshort{R2} and policy authorizes M.A to self-approve, \revshort{R2} becomes trusted and is delivered instead.

\section{History Event Model}
\begin{lstlisting}
pub struct HistoryEvent {
    pub event_id: [u8; 16],
    pub sequence: u64,
    pub file_id: [u8; 16],
    pub revision_id: Option<[u8; 32]>,
    pub previous_event_hash: [u8; 32],
    pub occurred_at: String,
    pub actor_identity: Option<[u8; 16]>,
    pub actor_label: Option<String>,
    pub topology_generation: Option<u64>,
    pub event_type: HistoryEventType,
    pub outcome: HistoryOutcome,
    pub details: EventDetails,
}
\end{lstlisting}

Suggested events include tracking start, checkout/check-in, revision signing, signature verification, countersignature, file requests, share attempts/delivery/accept/reject, unsigned fallback, access granted/denied, quorum use, parent approval, password/PIN use, share-link creation/redemption/revocation, bridge use, export, expiry, expired access, content destruction, policy change, rename, history extract/share, fork detection, tamper detection, and external modification detection.

\section{Secret-Safe Authentication History}
History may record that password or PIN authentication succeeded, failed, or locked out, but must never contain passwords, PINs, their hashes, KDF output, or other secret-derived material.
\begin{lstlisting}
PASSWORD_USED  actor: M.A  result: SUCCESS
PIN_USED       actor: M.A  result: FAILED  attempt_number: 3
\end{lstlisting}

\section{Quorum and Parent Approval History}
\begin{lstlisting}
QUORUM_USED
required_shares: 2
required_physical_devices: 2
shares_satisfied: 2
physical_devices_satisfied: 2
result: SUCCESS
\end{lstlisting}
History should record policy outcomes without storing share secrets.

\section{Direct Shares and Bearer Shares}
Direct recipient-addressed deliveries may record sender and recipient. Bearer links cannot always identify the redeemer and should use values such as \texttt{UNKNOWN\_BEARER} unless identity was independently authenticated.

\section{Expiration and Tombstones}
Tracked-file expiry should destroy content while retaining provenance:
\begin{lstlisting}
Before expiry:
KQTF
|-- R1(*@\revmark@*) payload
|-- R2(*@\revmark@*) payload
`-- history

After expiry:
KQTF
|-- payloads destroyed / unavailable
`-- history
    |-- FILE_EXPIRED
    `-- CONTENT_DESTROYED
\end{lstlisting}
The tombstone can still append \texttt{EXPIRED\_ACCESS\_ATTEMPT}. Expiry applies to all retained payload revisions so older signed content cannot bypass the expiry policy.

\section{History Integrity Chain}
\begin{lstlisting}
H0 = history genesis
H1 = SHA-256(KQ-FILE-HISTORY-v1 || H0 || canonical(Event1))
H2 = SHA-256(KQ-FILE-HISTORY-v1 || H1 || canonical(Event2))
history_root = Hn
\end{lstlisting}
Changing an older event invalidates later hashes.

Security-sensitive events may also carry domain-separated signatures:
\begin{lstlisting}
KQ-FILE-HISTORY-EVENT-v1
|| file_id
|| revision_id
|| sequence
|| previous_event_hash
|| event_hash
\end{lstlisting}

\section{Forks, Fast-Forward, and Merge}
If two holders extend a common base differently, preserve both branches and record a fork rather than silently overwriting one. If one side is merely ahead and the extension verifies, KeyQuorum may fast-forward. Never resolve by ``latest timestamp wins.''


\section{Conflict Detection and Human Resolution}

A history fork and a content conflict are related but distinct.

\begin{itemize}[leftmargin=1.5em]
  \item A \textbf{history fork} exists when two valid descendant revisions extend the same prior revision.
  \item A \textbf{content conflict} exists when those descendants contain incompatible changes that cannot be merged automatically without choosing between competing edits.
  \item A \textbf{policy conflict} exists when content may be mechanically mergeable but one or more branches do not satisfy the authorization, signature, topology, bridge, or trust policy required for the resulting revision.
\end{itemize}

KeyQuorum must never resolve a conflict by newest timestamp, by silently overwriting one branch, or merely by selecting the HCP with the lexicographically ``larger'' label. A human resolution path is determined from the file's own provenance history and HCP topology.

\subsection{Conflict Preservation}

When two current owners produce conflicting descendants, both revisions remain intact.

\begin{lstlisting}
               +--> R3A(*@\revmark@*)  M.A.2 edit
R2(*@\revmark@*) --------+
               `--> R3B(*@\revmark@*)  M.S.1 edit
\end{lstlisting}

If the changes do not overlap and the file type supports deterministic merging, KeyQuorum may construct a proposed merge revision. If the changes overlap, KeyQuorum records a conflict and selects a human reviewer using the escalation procedure below.

The conflicting revisions are never rewritten. A successful resolution creates a new multi-parent revision:

\begin{lstlisting}
R3A(*@\revmark@*) -----+
            +--> R4-MERGE(*@\revmark@*)
R3B(*@\revmark@*) -----+
\end{lstlisting}

The merge revision still passes through the normal signature, countersignature, bridge, topology, quorum, and policy evaluation before it can become trusted.


\subsection{Automatic Conflict Resolution Before Human Escalation}

Human escalation is the fallback, not the first action. When KeyQuorum detects two descendant heads from a common revision, it should first attempt a deterministic three-way merge using the nearest common ancestor revision as the merge base.

\begin{lstlisting}
                 R3A(*@\revmark@*)
                /
R2(*@\revmark@*)  <--- merge base
                \
                 R3B(*@\revmark@*)
\end{lstlisting}

The automatic resolver compares:

\begin{lstlisting}
base   = common ancestor content
left   = content of first conflicting head
right  = content of second conflicting head
\end{lstlisting}

The resolver should classify the result as:

\begin{lstlisting}
pub enum AutoMergeOutcome {
    FastForward,
    AlreadyEquivalent,
    CleanMerge,
    RequiresHuman,
    PolicyBlocked,
    UnsupportedContent,
}
\end{lstlisting}

The automatic path is intentionally conservative.

\subsubsection*{Automatic resolution rules}

KeyQuorum may resolve without human content review when one of the following is true:

\begin{enumerate}[leftmargin=1.8em]
  \item \textbf{Fast-forward}: one head is already an ancestor of the other and the descendant lineage verifies.
  \item \textbf{Equivalent result}: both heads independently produce the same resulting content commitment.
  \item \textbf{One-sided edit}: one head is unchanged from the merge base and the other head contains the only change.
  \item \textbf{Non-overlapping edits}: both heads changed the file, but their changed ranges do not overlap.
  \item \textbf{Identical overlapping edit}: both heads changed the same range to the same bytes.
\end{enumerate}

For an initial implementation, automatic content merging should be limited to content for which KeyQuorum has a deterministic merge strategy. UTF-8 text is the primary v1 candidate. Binary files and unsupported structured formats are preserved as divergent revisions and immediately enter the human-resolution path rather than being guessed.

For text, the resolver performs a true three-way merge against the common ancestor. It must preserve byte-level content deterministically and must not silently normalize unrelated data merely to make a merge succeed.

\begin{lstlisting}
base:
line 65: North
line 66: 100
line 67: South

left:
line 65: North
line 66: 125
line 67: South

right:
line 65: North
line 66: 100
line 67: South-East
\end{lstlisting}

Because the edits affect different ranges, KeyQuorum can construct:

\begin{lstlisting}
merged:
line 65: North
line 66: 125
line 67: South-East
\end{lstlisting}

The resulting revision has both heads as parents:

\begin{lstlisting}
R3A(*@\revmark@*) -----+
            +--> R4-MERGE(*@\revmark@*)
R3B(*@\revmark@*) -----+
\end{lstlisting}

\subsubsection*{Automatic merge does not imply automatic trust}

Automatic conflict resolution determines content, not authority. Signatures and policy evidence from the parent revisions do not automatically become signatures over the newly generated merge revision.

A machine-created merge therefore begins as a proposed revision:

\begin{lstlisting}
AUTO_MERGE_PROPOSED
parents:
  - R3A(*@\revmark@*)
  - R3B(*@\revmark@*)
result: CLEAN_MERGE

new_revision:
  R4-MERGE(*@\revmark@*)

trust_state:
  PENDING
\end{lstlisting}

The merged revision then passes through the ordinary KeyQuorum trust pipeline:

\begin{lstlisting}
automatic content merge
        |
        v
new multi-parent revision
        |
        v
human review of merge result
        |
        v
content signature required?
        |
        v
supervisor countersign required?
        |
        v
bridge / topology policy satisfied?
        |
        v
TRUSTED
        |
        v
finalization signature
        |
        v
FINAL
\end{lstlisting}

By default, a deterministic auto-merge must not inherit a parent's content signature. The new revision contains new bytes and a new revision ID, so the applicable file policy must decide who signs or countersigns it.


\subsubsection*{Native CLI Human Review and Line Provenance}

Human review should be part of the normal merge workflow, not a separate long-form command requiring many parameters. When a merge produces a reviewable result, or when automatic resolution requires a person, the native CLI should enter a compact interactive review view when attached to an interactive terminal.

Typical usage remains simple:

\begin{lstlisting}
keyquorum file merge decision.txt
\end{lstlisting}

If the merge is clean but requires review, the command opens the review interface automatically. If a pending review already exists, the user can reopen it with:

\begin{lstlisting}
keyquorum file review decision.txt
\end{lstlisting}

The user should not need to provide revision IDs, parent hashes, conflict IDs, or HCP labels for the common review path. KeyQuorum derives those values from the pending merge/review state. Explicit low-level flags may remain available for scripting or diagnostics, but they are not the normal human workflow.

\paragraph{Changed lines first.}

The review view should place changed lines at the top so the reviewer sees the substantive edits before unchanged context. Recommended order:

\begin{enumerate}[leftmargin=1.8em]
  \item changed ranges;
  \item conflict or merge explanation;
  \item signer/reviewer/finalization state;
  \item surrounding unchanged context; and
  \item optional full-file view.
\end{enumerate}

The same color convention is used throughout:

\begin{itemize}[leftmargin=1.5em]
  \item \diffadd{Green background + dark green text}: added or replacement content accepted into the candidate merge.
  \item \diffremove{Red background + dark red text}: removed or superseded content.
  \item \diffsame{White background + dark neutral text}: unchanged context.
\end{itemize}

Conceptual terminal layout:

\begin{lstlisting}
decision.txt -- merge review

CHANGED LINES

- 66 | approve project on Friday
+ 66 | approve project on Monday

+ 84 | notify accounting after approval

MERGE
  base:    a41c9e7...
  left:    0f6a3c1...
  right:   7be219d...
  result:  9c44d72...

STATUS
  merge:   CLEAN
  review:  PENDING
  trust:   PENDING
  final:   NO
\end{lstlisting}

The actual terminal renderer applies the red/green/white backgrounds described above rather than relying on the ASCII \texttt{+}/\texttt{-} markers alone.

\paragraph{Mouse-hover provenance.}

Each changed line or changed range should carry provenance metadata identifying the revision and HCP that introduced it. In the interactive CLI review, moving the mouse over a changed line reveals a small information panel.

For example, hovering over the green replacement for line 66 may show:

\begin{lstlisting}
Changed by: M.A.2
Revision:   0f6a3c1...
When:       2026-10-02T14:32:05Z
Parent:     a41c9e7...
Source:     LEFT BRANCH
\end{lstlisting}

Hovering over a red removed line may show:

\begin{lstlisting}
Previous author: M.A
Revision:        a41c9e7...
Superseded by:   M.A.2
\end{lstlisting}

This is presentation over stored revision provenance. The TUI must not infer authorship from terminal text. The merge/revision layer supplies the provenance record and the terminal renderer only displays it.

A useful internal view model is:

\begin{lstlisting}
pub struct ChangedRangeView {
    pub start_line: u32,
    pub end_line: u32,
    pub kind: ChangeKind,
    pub revision_id: [u8; 32],
    pub author_identity: Option<[u8; 16]>,
    pub author_hcp_label: String,
    pub changed_at_utc: String,
    pub source_parent: MergeParent,
}
\end{lstlisting}

For a line composed from both parents, the view may contain multiple provenance contributors rather than fabricating one author.

\paragraph{Terminal implementation.}

The native interactive renderer may use \texttt{ratatui} with \texttt{crossterm} mouse capture. Hover-sensitive rows maintain screen hit regions associated with \texttt{ChangedRangeView}. Mouse movement selects the region and renders the provenance panel.

This UI layer contains no merge, trust, signature, or HCP policy logic. It consumes the same merge/revision objects used by the non-interactive CLI path.

Keyboard access must mirror mouse access:

\begin{lstlisting}
j / k       move down / up one changed line
h / l       move to previous / next review pane
gg          jump to first changed line
G           jump to last changed line
Ctrl-d      move down half a page
Ctrl-u      move up half a page
Enter       inspect provenance for selected line
za          toggle surrounding unchanged context
/           search within the review
n / N       next / previous search result
i           open selected merge result for manual editing
:accept     accept the merge result
:reject     reject the merge result
:sign       sign the accepted revision
:finalize   add the finalization signature
:q          leave review
Esc         leave command/search mode
\end{lstlisting}

The review interface should follow Vim-style navigation rather than arrow-key-first navigation. Single-key movement uses familiar Vim motions, while security-sensitive state transitions use explicit command-mode verbs. This keeps navigation fast without making acceptance, rejection, signing, or finalization easy to trigger accidentally.

Arrow keys may remain as optional compatibility aliases for \texttt{j}/\texttt{k}, but the documented interface is Vim-oriented. Destructive or trust-changing operations should require an explicit colon command and confirmation when policy requires it.


If mouse reporting is unavailable, selecting a changed line with the keyboard shows the same provenance panel.

\paragraph{Review actions.}

The normal human review should therefore feel like:

\begin{lstlisting}
keyquorum file merge decision.txt

    -> KeyQuorum detects branch merge
    -> auto-merge runs
    -> changed lines appear first
    -> reviewer hovers/selects lines to inspect authorship
    -> reviewer accepts, rejects, or edits
    -> applicable signatures/countersignatures run
    -> authorized finalizer signs
    -> revision becomes FINAL
\end{lstlisting}

This workflow keeps merge review concise for humans while retaining lower-level scriptable commands for automation.


\subsubsection*{Conflicts that automatic resolution must not decide}

Automatic merge must stop and escalate when:

\begin{itemize}[leftmargin=1.5em]
  \item both heads modify the same range differently;
  \item a binary or unsupported format has no deterministic merge strategy;
  \item the revision DAG or history chain fails verification;
  \item a parent revision is tampered, expired, revoked, or otherwise ineligible;
  \item topology or policy evidence disagrees in a way that requires an authority decision;
  \item private bridge policy prevents enough information from being disclosed to make the merge decision safely;
  \item merging would require choosing which HCP's security decision overrides another HCP's decision; or
  \item file policy explicitly disables automatic merging.
\end{itemize}

In particular, KeyQuorum must never use HCP seniority to pick one conflicting line over another automatically. Seniority, prior ownership, sector prominence, and common-ancestor escalation are used to select the \emph{human resolver} after deterministic content resolution has failed.

\subsubsection*{Automatic merge events}

Recommended events are:

\begin{lstlisting}
AUTO_MERGE_ATTEMPTED
AUTO_MERGE_FAST_FORWARD
AUTO_MERGE_EQUIVALENT
AUTO_MERGE_CLEAN
AUTO_MERGE_BLOCKED
AUTO_MERGE_REQUIRES_HUMAN
\end{lstlisting}

Example:

\begin{lstlisting}
AUTO_MERGE_ATTEMPTED
base_revision: R2(*@\revmark@*)
left_revision: R3A(*@\revmark@*)
right_revision: R3B(*@\revmark@*)

AUTO_MERGE_CLEAN
strategy: THREE_WAY_TEXT
new_revision: R4-MERGE(*@\revmark@*)
trust_state: PENDING

POLICY_DECISION
revision: R4-MERGE(*@\revmark@*)
result: PENDING
reason: CONTENT_SIGNATURE_REQUIRED
\end{lstlisting}

If the automatic resolver reaches \texttt{RequiresHuman}, KeyQuorum then runs the human resolver-selection algorithm below.


\subsection{Owner History}

For conflict resolution, an \textbf{owner} means an HCP that appears in the tracked file's provenance as an authorized holder or participant in the file's custody or revision lineage. Ownership for this purpose is derived from authenticated file history, not merely from possession of a local copy.

After automatic resolution returns \texttt{RequiresHuman}, the resolver-selection algorithm considers:
\begin{enumerate}[leftmargin=1.8em]
  \item the current conflicting owners;
  \item prior owners recorded before those conflicting owners acquired or modified the file;
  \item HCP ancestry and depth;
  \item historical sector representation;
  \item bridge relationships between sectors; and
  \item the least common senior HCP that can authoritatively connect the unresolved branches.
\end{enumerate}

\subsection{Rule 1: Prefer the Most Recent Prior Neutral Owner}

The first choice is the most recent prior owner in the file's authenticated history who is \textbf{not} one of the current owners whose revisions created the conflict.

\begin{lstlisting}
decision.txt ownership / custody history:

M.A
 |
 +--> M.A.2
 |
 `--> M.S.1

Current conflicting owners:
  M.A.2
  M.S.1

Prior neutral owner:
  M.A
\end{lstlisting}

Assume \texttt{M.A.2} and \texttt{M.S.1} are connected through an authorized bridge and both edit line 66 differently. Since \texttt{M.A} was a prior owner and is not one of the two conflicting editors, KeyQuorum assigns resolution review to \texttt{M.A}.

\begin{lstlisting}
CONTENT_CONFLICT_DETECTED
file: decision.txt
revision_a: R3A(*@\revmark@*)
actor_a: M.A.2
revision_b: R3B(*@\revmark@*)
actor_b: M.S.1
conflict: line 66

CONFLICT_REVIEW_ASSIGNED
reviewer: M.A
selection_rule: PRIOR_NEUTRAL_OWNER
\end{lstlisting}

The reviewer does not modify either parent revision. The reviewer approves one branch, rejects one branch, or produces/approves a new merge revision with both conflicting revisions as parents.

\subsection{Rule 2: If Prior Ownership Does Not Produce a Unique Reviewer, Prefer the Most Senior Historical Owner}

If Rule 1 cannot identify one valid reviewer, KeyQuorum evaluates the HCPs represented in the file's authenticated ownership history by topology seniority.

Seniority is determined structurally:
\begin{itemize}[leftmargin=1.5em]
  \item An ancestor is senior to its descendant.
  \item A shallower HCP depth is senior to a deeper HCP depth within the same connected hierarchy.
  \item \texttt{M} is senior to \texttt{M.A} and \texttt{M.S}; \texttt{M.S} is senior to \texttt{M.S.1}.
  \item Sibling branches at the same depth, such as \texttt{M.A} and \texttt{M.S}, are equal in seniority unless another rule below resolves the tie.
\end{itemize}

This is a topology rule, not a job-title guess.

\subsection{Rule 3: If Seniority Ties Across Sectors, Use Historical Sector Prominence}

If the most senior eligible owners are tied because they belong to different branches at the same HCP depth, KeyQuorum counts historical owners by sector for this file.

\begin{lstlisting}
Historical owners represented:

Accounting sector:
  M.A
  M.A.1
  M.A.2
  total = 3

Software sector:
  M.S
  M.S.1
  total = 2

Most senior tied owners:
  M.A
  M.S
\end{lstlisting}

The Accounting sector is more prominent in this file's authenticated ownership history. KeyQuorum therefore resolves the sector tie in favor of the Accounting branch and then reapplies seniority within that sector. Because \texttt{M.A} is present and is senior to \texttt{M.A.1} and \texttt{M.A.2}, \texttt{M.A} becomes the human resolver.

Sector prominence is a tie-breaker only. It does not allow a lower descendant to outrank an available ancestor in the same sector.

\subsection{Rule 4: If Prior Owner, Seniority, and Sector Prominence All Fail, Escalate to a Common Senior Ancestor}

If the preceding rules still produce a tie, KeyQuorum escalates to the closest HCP ancestor that can authoritatively connect the unresolved sectors without depending on the conflicting bridge itself.

\begin{lstlisting}
Historical owners:

M.A
M.A.1
M.S
M.S.1

sector totals:
    M.A branch = 2
    M.S branch = 2

most senior sector owners:
    M.A
    M.S

seniority:
    tied

sector prominence:
    tied
\end{lstlisting}

The closest common senior ancestor is \texttt{M}:

\begin{lstlisting}
M
|-- M.A
|   `-- M.A.1
`-- M.S
    `-- M.S.1
\end{lstlisting}

Therefore \texttt{M} receives final conflict review authority.

\begin{lstlisting}
CONFLICT_REVIEW_ESCALATED
from:
  M.A
  M.S
to:
  M
selection_rule:
  LOWEST_COMMON_SENIOR_ANCESTOR
\end{lstlisting}

This is equivalent to using the least common ancestor of the unresolved HCP branches, subject to file policy and the requirement that the resulting HCP actually has authority to resolve the file conflict.

\subsection{Bridge-Assisted Escalation}

Bridges do not automatically determine a winner. Their role in conflict handling is to make cross-branch relationships explicit and provide an authenticated path for review, notification, countersigning, and escalation.

For the final escalation case, a bridge between two sectors can help KeyQuorum identify the intended cross-branch supervisors or bridge-visible ancestry and route the conflict review earlier.

\begin{lstlisting}
BRIDGE_USED
operation: CONFLICT_ESCALATION
from_branch: M.A
to_branch: M.S
result: REVIEW_PATH_ESTABLISHED
\end{lstlisting}

A private bridge must preserve its existing privacy rules. It may satisfy policy internally without exposing bridge membership or identifying details in transferable file history.

\subsection{Resolver Selection Algorithm}

\begin{lstlisting}
conflicting_owners =
    owners that produced conflicting heads

1. prior_candidates =
       authenticated prior owners
       - conflicting_owners

2. if exactly one valid most-recent
   prior neutral owner exists:
       resolver = that owner
       stop

3. senior_candidates =
       most senior eligible HCPs
       in ownership history

4. if exactly one senior candidate exists:
       resolver = senior candidate
       stop

5. group owners by HCP sector
   compare historical sector prominence

6. if one sector is dominant:
       resolver =
           most senior eligible HCP
           in dominant sector
       stop

7. resolver =
       nearest common senior ancestor
       of unresolved branches

8. if no authorized common ancestor exists:
       mark conflict UNRESOLVED
       require explicit root/admin policy
\end{lstlisting}

The algorithm must operate on authenticated provenance records and topology generations relevant to the file history. It must not rely on mutable UI labels or current organizational assumptions alone.

\subsection{Conflict Events}

Recommended events include:

\begin{lstlisting}
HISTORY_FORK_DETECTED
CONTENT_CONFLICT_DETECTED
POLICY_CONFLICT_DETECTED
CONFLICT_REVIEW_ASSIGNED
CONFLICT_REVIEW_ESCALATED
MERGE_PROPOSED
CONFLICT_RESOLVED
MERGE_REJECTED
\end{lstlisting}

A full resolution may look like:

\begin{lstlisting}
CONTENT_CONFLICT_DETECTED
revision_a: R3A(*@\revmark@*)
revision_b: R3B(*@\revmark@*)
common_parent: R2(*@\revmark@*)

CONFLICT_REVIEW_ASSIGNED
reviewer: M.A
selection_rule: PRIOR_NEUTRAL_OWNER

MERGE_PROPOSED
actor: M.A
parents:
  - R3A(*@\revmark@*)
  - R3B(*@\revmark@*)

CONFLICT_RESOLVED
actor: M.A
resolution: MANUAL
new_revision: R4-MERGE(*@\revmark@*)

REVISION_SIGNED
actor: M.A
revision: R4-MERGE(*@\revmark@*)

POLICY_DECISION
revision: R4-MERGE(*@\revmark@*)
result: TRUSTED
\end{lstlisting}

\subsection{Conflict Resolution Invariants}

\begin{enumerate}[leftmargin=1.8em]
  \item Conflicting parent revisions are immutable and remain in the DAG.
  \item A human resolution creates a new revision; it never rewrites either parent.
  \item Current conflicting owners do not automatically resolve their own collision.
  \item The most recent prior neutral owner is preferred when uniquely identifiable and authorized.
  \item HCP ancestry determines seniority; names and job titles do not.
  \item Sector prominence is used only to break otherwise equal cross-sector seniority.
  \item If sector prominence also ties, resolution escalates to the nearest authorized common senior ancestor.
  \item Bridges assist routing and cross-branch authorization but do not automatically decide the winner.
  \item Private bridge metadata remains private.
  \item The resolver's decision still requires the normal revision trust policy before a merge becomes trusted.
  \item If no authorized resolver can be found, the conflict remains explicitly unresolved rather than being guessed.
\end{enumerate}


\section{Timestamp Model}
Use UTC timestamps, but do not treat client clocks as authoritative ordering. Event sequence and hash linkage establish order. Relay-observed time may optionally supplement actor-reported time.

\section{SQLite as an Index, Not the Authority}
SQLite may contain rebuildable indexes such as \texttt{tracked\_files} and \texttt{tracked\_history\_index}. If the cache is lost, KeyQuorum should rebuild it by verifying KQTF.

\section{Suggested Module Layout}
\begin{lstlisting}
src/file_history.rs
src/file_history/
|-- codec.rs
|-- container.rs
|-- event.rs
|-- revision.rs
|-- topology.rs
|-- policy.rs
|-- merge.rs
|-- verify.rs
`-- tests.rs
\end{lstlisting}

Suggested API includes \texttt{start\_tracking}, \texttt{append\_event}, \texttt{check\_in\_revision}, \texttt{sign\_revision}, \texttt{countersign\_revision}, \texttt{evaluate\_revision\_trust}, \texttt{select\_shareable\_revision}, \texttt{verify\_history}, \texttt{expire\_payload}, \texttt{extract\_history}, and \texttt{merge\_received\_history}.

\section{File Delivery and History Transport}
Existing KQPB file delivery should bind the delivery ID, sender, recipient, file ID, revision ID, history root, logical filename, tracked payload, delivery-policy result, and content-proof descriptor. The receive path verifies transport signature, KQTF history, revision proof, topology/policy, and then merges history and records acceptance.

Extracted history should be structured protocol data, not an ordinary user file:
\begin{lstlisting}
pub struct HistorySnapshot {
    pub file_id: [u8; 16],
    pub through_sequence: u64,
    pub history_root: [u8; 32],
    pub events: Vec<HistoryEvent>,
}
\end{lstlisting}
KQPB may gain dedicated history kinds and KQXB may gain a history-snapshot subtype.

\section{Privacy and Attribution}
History must never expose private keys, raw API keys, bridge shared secrets, passwords, PINs, or bearer tokens. It must never claim more attribution than KeyQuorum can establish. Unknown actors remain \texttt{UNKNOWN}; unauthenticated bearer users remain \texttt{UNKNOWN\_BEARER}.

\section{Core Invariants}
\begin{enumerate}[leftmargin=1.8em]
  \item Newer does not imply trusted.
  \item Trust requires the applicable cryptographic and topology policy.
  \item History travels with KQTF; SQLite is only cache/index.
  \item File identity is independent of filename/path.
  \item Revision identity is independent of file identity.
  \item Authority is evaluated against the topology generation applicable to the event.
  \item Supervisor approval is an explicit cryptographic action.
  \item Private bridge details do not leak through portable history.
  \item Authentication secrets never appear in history.
  \item Expired content may be destroyed while provenance remains.
  \item Expired tracked files can still record attempted access.
  \item Forks are never silently overwritten.
  \item Attribution is never fabricated.
  \item Transport authentication and content trust are separate.
\end{enumerate}

\section{Alignment With the Current KeyQuorum Repository}

The history architecture should be implemented as an upgrade to the current KeyQuorum execution paths, not as a parallel copy of existing security logic. The current repository already contains the primitives for ancestry, signatures, bridge membership, quorum reconstruction, password/PIN gates, expiry, sharing, device identity, recipient-sealed transport, and lab activity tracing. File History should compose those primitives and add provenance state around their existing decisions.

\subsection{Implementation Rule: Upgrade Shared Code; Do Not Duplicate Security Methods}

The implementation should follow one rule throughout:

\begin{quote}
\textbf{If KeyQuorum already has the method that decides or verifies something, File History records and reuses that result; it does not reimplement the decision in \texttt{file\_history}.}
\end{quote}

Examples:

\begin{itemize}[leftmargin=1.5em]
  \item Do not create a second Ed25519 verifier in \texttt{file\_history}; call \texttt{signing::verify\_signature}.
  \item Do not create a second parent-label parser; promote/reuse the existing topology helpers.
  \item Do not create a second KQPB serializer; extend \texttt{envelope.rs}.
  \item Do not create a second file-delivery protocol; extend \texttt{file\_delivery.rs}.
  \item Do not create a second quorum evaluator; let \texttt{quorum.rs} make the access decision and emit history metadata around that result.
  \item Do not put a second conflict engine in the browser lab; the lab invokes the same Rust/CLI implementation compiled to WASM.
\end{itemize}

The new \texttt{file\_history} subsystem owns revision DAGs, KQTF encoding, history chaining, conflict classification, merge construction, and provenance presentation. Existing modules continue to own their existing security domains.

\subsection{Topology and Authority: \texttt{private\_bridge.rs} and \texttt{authority.rs}}

The repository already has reusable HCP topology primitives:

\begin{itemize}[leftmargin=1.5em]
  \item \texttt{private\_bridge::parent\_node\_label}: one-step dotted ancestry such as \texttt{M.S.2 -> M.S}.
  \item \texttt{private\_bridge::is\_ancestor\_or\_self}: segment-aware ancestry recognition.
  \item \texttt{authority::restructure\_countersigner}: current parent-countersign selection.
  \item \texttt{authority::is\_routine\_employee\_reissue}: delegated direct-descendant policy.
  \item \texttt{authority::require\_unlock\_approval}: current opt-in parent approval verification for file unlocks.
\end{itemize}

File History should not duplicate those rules in \texttt{src/file\_history/topology.rs}. Instead, the codebase should upgrade the existing topology layer into a single shared API used by bridges, organization updates, transfers, unlock approval, revision trust, and conflict escalation.

A practical refactor is:

\begin{lstlisting}
authority / shared topology API:
    parent_node_label(...)
    is_ancestor_or_self(...)
    ancestry_distance(...)
    relationship(...)
    lowest_common_ancestor(...)
    direct_parent(...)
\end{lstlisting}

Existing callers in \texttt{private\_bridge.rs}, \texttt{transfer.rs}, \texttt{org\_update}, and the lab should then call or re-export the shared implementation. During migration, compatibility wrappers may remain, but there must be only one underlying algorithm.

The conflict resolver uses these shared functions for seniority, sector ancestry, prior-owner eligibility, and the final common-ancestor escalation. Revision trust uses the same functions for scope-owner, descendant, ancestor, and cross-branch classification.

\subsection{Bridge State: \texttt{private\_bridge.rs}}

The current bridge subsystem already owns bridge generation, membership, supervisor metadata, private shared signing material, notification labels, and bridge lifecycle events. File History should query that existing bridge state instead of maintaining a second bridge graph.

For non-private history disclosure, the history layer may record that an authorized bridge participated in revision approval, sharing, or conflict escalation. For private bridges, the bridge module remains the authority on what may be disclosed; File History records only the policy outcome that is safe to expose.

Bridge-assisted conflict escalation therefore asks the bridge subsystem for the applicable relationship/authorization result and then records:

\begin{lstlisting}
BRIDGE_USED
operation: CONFLICT_ESCALATION
result: REVIEW_PATH_ESTABLISHED
\end{lstlisting}

when disclosure policy permits.

\subsection{Signatures: \texttt{signing.rs}}

The current signing module already provides:

\begin{itemize}[leftmargin=1.5em]
  \item strict Ed25519 verification through \texttt{verify\_signature};
  \item signing through \texttt{sign};
  \item bridge-plus-personal dual signatures through \texttt{sign\_with\_bridge}; and
  \item bridge signature verification through \texttt{verify\_bridge\_signature}.
\end{itemize}

The history upgrade should add only new domain-separated preimage constructors and artifact codecs where required, for example:

\begin{lstlisting}
KQ-FILE-REVISION-v1
KQ-FILE-COUNTERSIGN-v1
KQ-FILE-HISTORY-EVENT-v1
\end{lstlisting}

Those preimages still call the existing \texttt{sign} and \texttt{verify\_signature} functions. File History must not add its own Ed25519 implementation.

The same rule applies to bridge signatures: revision policy references/verifies existing KQBS semantics rather than inventing a second bridge-signature format.

\subsection{Quorum and Physical Custody: \texttt{quorum.rs}, \texttt{key\_tree.rs}, and \texttt{device.rs}}

The current quorum path already reconstructs presented key-tree shares, tracks the devices that contributed, enforces physical-device policy, invokes \texttt{authority::require\_unlock\_approval}, decrypts the protected file, and writes \texttt{unlock\_events} for success or failure.

File History should wrap this path with structured provenance rather than re-evaluate quorum independently.

Conceptually:

\begin{lstlisting}
existing quorum unlock
        |
        +--> existing key-tree reconstruction
        +--> existing physical-device evaluation
        +--> existing parent approval
        +--> existing AES-GCM decrypt
        |
        v
history adapter records outcome
\end{lstlisting}

For tracked files, the long-term upgrade should move from an audit row that disappears with the file to a structured KQTF event while retaining a rebuildable SQLite index if useful. Existing \texttt{unlock\_events} may be kept during migration for compatibility, but the authoritative tracked history becomes the KQTF event chain.

\subsection{Password Protection and Expiry: \texttt{locked\_files.rs}}

The current password-file code already owns password KDF/decryption, expiry parsing, future-expiry validation, expiration checks, and ciphertext destruction.

The history subsystem should not reproduce any of those methods. Instead:

\begin{itemize}[leftmargin=1.5em]
  \item call the existing unlock function;
  \item record only success/failure and actor context, never the password;
  \item reuse the existing UTC expiry parser and validation logic; and
  \item upgrade tracked-file expiry so content destruction leaves a KQTF tombstone/history rather than deleting the authoritative provenance.
\end{itemize}

For non-tracked legacy files, current behavior may remain unchanged.

\subsection{PIN Gates: \texttt{pin.rs}}

The current PIN module already defines whether verification is required, performs Argon2id-backed PIN verification, counts failed attempts, locks after the configured threshold, and manages temporary unlock windows.

File History should leave this logic generic. The operation that calls \texttt{pin::verify\_pin} records the result:

\begin{lstlisting}
PIN_USED
result: SUCCESS | FAILED | LOCKED
\end{lstlisting}

No PIN checking code belongs in \texttt{file\_history}.

\subsection{Share Links: \texttt{sharing.rs}}

The current sharing module already creates one-time bearer tokens, stores only token hashes, enforces expiration, max-use counts, revocation, and atomically increments redemption use counts.

The history upgrade should instrument the existing create/redeem/revoke boundaries:

\begin{lstlisting}
create_file_share(...)
    -> SHARE_LINK_CREATED

redeem_file_share(...)
    -> SHARE_LINK_REDEEMED
       or ACCESS_DENIED

revoke_file_share(...)
    -> SHARE_LINK_REVOKED
\end{lstlisting}

The existing token and concurrency logic remains unchanged. If the bearer is not otherwise authenticated, the history actor stays \texttt{UNKNOWN\_BEARER}.

\subsection{Authenticated File Delivery: \texttt{file\_delivery.rs}}

The current file-delivery protocol already supplies a random delivery ID, sender/recipient labels, recipient-bound encryption, a sender transport signature, a content hash, and a signed recipient acknowledgement.

Do not add a second ``tracked delivery'' implementation. Upgrade the existing \texttt{Outgoing}, \texttt{FileLetter}, \texttt{DeliveryAck}, and their preimages/codecs to optionally carry:

\begin{lstlisting}
file_id
revision_id
history_root
generated_label
optional_user_label
delivery_policy_result
content_proof_descriptor
KQTF payload
\end{lstlisting}

Legacy/untracked delivery can remain a versioned compatibility path. Tracked delivery uses the same \texttt{seal\_letter}, \texttt{open\_letter}, \texttt{seal\_ack}, and \texttt{open\_ack} workflow after its wire format is upgraded/versioned.

This keeps transport authentication and content trust separate while avoiding duplicate protocols.

\subsection{Envelope and Codec Reuse: \texttt{envelope.rs}}

The current envelope module already centralizes KQPB/KQXB framing, X25519 recipient sealing, weak-public-key rejection, parsing, and the length-prefixed codec helpers used by signed preimages.

History transport should extend this module with new message kinds rather than define another envelope layer.

For example:

\begin{lstlisting}
KIND_FILE_HISTORY
KIND_FILE_HISTORY_ACK
\end{lstlisting}

History payload encoders should reuse \texttt{push\_len\_prefixed}, \texttt{push\_len\_prefixed\_u32}, \texttt{take\_array}, \texttt{take\_len\_prefixed}, and related canonical framing helpers so the new wire format follows the same bounds and parsing model already used by KeyQuorum.

\subsection{Stable Identity and Device Movement: \texttt{transfer.rs}}

The current transfer subsystem already separates stable 16-byte cryptographic identity from dotted HCP placement and device-local possession, including \texttt{active}/\texttt{ghost}, COPY/MOVE semantics, descendant modes, and transaction/audit state.

File History should reference existing stable HCP identity IDs for attribution where available instead of inventing a file-specific person identifier.

File ownership history remains file-specific and should not be confused with \texttt{key\_provenance}; however, the same stable identity is used to say that the \emph{same principal} authored, received, transferred, or later resolved a revision even if device custody changes.

\subsection{SQLite: \texttt{db/schema.sql} and \texttt{db/mod.rs}}

The database remains a cache/index and migration surface, not the provenance authority for tracked files.

Schema upgrades should add rebuildable indexes for tracked files, revisions, and history events. Existing migration machinery in \texttt{db/mod.rs} should be extended to add or rebuild those tables instead of creating an independent migration mechanism inside \texttt{file\_history}.

Likewise, transaction boundaries should reuse the existing database transaction helpers where applicable.

\subsection{Automatic Merge: New Logic, Existing Trust Primitives}

Automatic three-way content merge is genuinely new behavior, so it belongs in \texttt{src/file\_history/merge.rs}. Even there, only the content-diff/merge algorithm is new.

The merge module must call existing/shared primitives for everything else:

\begin{lstlisting}
revision ancestry     -> file_history/revision.rs
HCP relationship      -> shared authority/topology API
bridge authorization  -> private_bridge.rs
signature verification-> signing.rs
trust decision        -> file_history/policy.rs
wire transport        -> file_delivery.rs + envelope.rs
\end{lstlisting}

The merge engine never reimplements signatures, ancestry, bridge membership, or delivery.

\subsection{Lab Architecture: Reuse the Existing Activity Pipeline}

The current browser lab already has an \textbf{Activity} tab in \texttt{lab/src/App.tsx}. Its \texttt{ActivityPanel.tsx} renders the latest access trace and the full \texttt{snapshot.activity} log. The Rust side, \texttt{src/lab/view.rs}, defines \texttt{ActivityView}, and \texttt{src/lab/state.rs} is the single lab state used by GUI buttons, the terminal, and the WASM facade.

That is the correct home for File History in the lab. Do \textbf{not} add a separate History tab or a second React-side event store.

Upgrade the existing activity projection instead:

\begin{lstlisting}
ActivityView
    existing:
        seq
        actor
        kind
        outcome
        title
        trace
        command

    add optional tracked-file context:
        file_id
        revision_id
        generated_label
        user_label
        history_event_type
        history_root
        conflict_id
        parent_revision_ids
        review_state
        finalization_state
\end{lstlisting}

The Activity page can then provide:

\begin{itemize}[leftmargin=1.5em]
  \item the existing latest-operation trace;
  \item the existing full activity log;
  \item filters for file, actor, event type, and trust outcome;
  \item expandable tracked-file history entries;
  \item revision IDs/generated labels;
  \item conflict and auto-merge events;
  \item resolver assignment/escalation events; and
  \item links or expandable views for a selected file's revision DAG.
\end{itemize}

The browser UI remains a renderer. Per the current lab architecture, \texttt{src/lab/state.rs} runs real \texttt{keyquorum}/\texttt{keyquorum-device} commands through the lab VM, and the VM itself does not decide authorization. The history implementation must follow the same rule: the real Rust/CLI path performs tracking, merge, signature, trust, and conflict decisions; the lab receives those results through the existing snapshot/activity projection.

The TypeScript \texttt{lab/src/api/types.ts} should continue to mirror \texttt{src/lab/view.rs}; do not define a second independent history model in TypeScript.

\subsection{Upgrade Path and Compatibility}

To avoid method duplication during implementation:

\begin{enumerate}[leftmargin=1.8em]
  \item Add KQTF/revision/history primitives in \texttt{file\_history}.
  \item Promote shared topology helpers once, then migrate existing callers.
  \item Add revision/history domain constructors to \texttt{signing.rs}; continue using the existing signing engine.
  \item Version/extend existing file-delivery and envelope codecs rather than creating parallel transports.
  \item Add optional history emission hooks/context to existing file operations.
  \item Keep legacy untracked file operations working without forcing them into KQTF.
  \item Add rebuildable database indexes and migrations through the existing DB migration path.
  \item Extend the lab's existing \texttt{ActivityView}/\texttt{ActivityPanel}; do not build a second activity/history pipeline.
  \item Remove compatibility wrappers only after all existing callers use the shared methods.
\end{enumerate}

The result is one security implementation with additional provenance, not two competing implementations that can drift apart.



\section{Relationship to Git}
KeyQuorum File History has a meaningful resemblance to Git, but it should not be described as ``Git embedded in a file.'' Git is a distributed version-control system optimized for source trees and developer collaboration. KeyQuorum File History is a cryptographic provenance and policy-enforcement system optimized for secure documents, HCP authority, hardware custody, and controlled sharing.

The resemblance is strongest at the data-structure level: immutable revisions, parent relationships, branch/fork detection, merge semantics, content commitments, and signed checkpoints. The differences are strongest at the security layer: HCP topology, supervisor countersigning, quorum, device custody, expiry, password/PIN outcomes, recipient-bound delivery, and policy-controlled revision selection.

\subsection{Concept Mapping}
\begin{longtable}{>{\raggedright\arraybackslash}p{0.22\textwidth} >{\raggedright\arraybackslash}p{0.34\textwidth} >{\raggedright\arraybackslash}p{0.34\textwidth}}
\toprule
\textbf{Git Concept} & \textbf{KeyQuorum Analogue} & \textbf{Important Difference} \\
\midrule
\endhead
Repository & KQTF tracked-file container & KQTF is normally scoped to one tracked document and carries security provenance. \\
Commit & Revision checkpoint & A KeyQuorum revision may remain untrusted until HCP/signature/countersignature/bridge policy succeeds. \\
Commit hash & Cryptographic revision ID & The KeyQuorum revision ID is a SHA-256 commitment to the canonical revision object, including its parent revision ID(s), while protected content uses a separate content commitment to avoid blindly exposing plaintext fingerprints. \\
Parent commit & Parent revision & KeyQuorum can use one parent for normal edits or multiple parents for explicit merges. \\
Branch & Divergent revision lineage & Divergence may represent a security/provenance conflict requiring policy resolution. \\
Merge commit & Approved merged revision & The merge may require supervisor, bridge, scope-owner, or root approval. \\
Signed commit/tag & Content signature/trusted checkpoint & Signature validity is only one input to KeyQuorum trust. \\
Author & Revision author HCP & Identity is cryptographic and topology-aware, not a free-form name/email field. \\
Committer & Approver/countersigner & Approval can be cryptographically distinct from authorship. \\
Reflog & Local audit/index cache & KeyQuorum's authoritative history travels with the file. \\
Remote & Recipient/peer/relay & KeyQuorum uses recipient-sealed and policy-aware transport. \\
Fast-forward & Verified history extension & Allowed only after lineage, proofs, topology, and policy validate. \\
Divergence & History fork & Preserve and explicitly resolve; never silently overwrite. \\
Tag & Named trusted checkpoint & Trust remains policy- and signature-dependent. \\
.git directory & Embedded KQTF metadata/proofs/history & Provenance is intended to remain attached to the tracked document. \\
\bottomrule
\end{longtable}

\subsection{The Strongest Similarity: Immutable Revision Lineage}
Git's central structural idea is an immutable graph where each commit identifies its parent commit or commits. KeyQuorum adopts the same structural rule: every production revision ID cryptographically commits to its parent revision ID or IDs.

\begin{lstlisting}
a41c9e7... -> bc2810f... -> f93ad44...
\end{lstlisting}

For concurrent edits:

\begin{lstlisting}
                +--> 0f6a3c1...
a41c9e7... -----+
                `--> 7be219d...
\end{lstlisting}

An explicit merge can produce:

\begin{lstlisting}
0f6a3c1... ----+
               +--> 9c44d72...
7be219d... ----+
\end{lstlisting}

Human-readable generated labels and optional user labels accompany these hashes but do not replace them. Shorthand such as \revshort{R1}/\revshort{R2A} is permitted only in explanatory material. This is structurally close to Git's commit DAG and is preferable to assuming all distributed editing is linear.

\subsection{Recommended Refinement: Revision DAG Versus Event Chain}
The audit history and the content lineage are related but different structures.
\begin{enumerate}[leftmargin=1.8em]
  \item \textbf{Revision DAG}: which content revision descended from which content revision(s).
  \item \textbf{History/Event Chain}: the authenticated sequence in which security-relevant events were recorded.
\end{enumerate}

\begin{lstlisting}
REVISION DAG
                +--> 0f6a3c1...
a41c9e7... -----+
                `--> 7be219d...

EVENT CHAIN
E1 -> E2 -> E3 -> E4 -> E5
\end{lstlisting}
The event chain may record \texttt{HISTORY\_FORK\_DETECTED}, while the revision DAG preserves both branches. This separation is one of the most useful ideas KeyQuorum can borrow from Git.

\subsection{Git Commit Versus KeyQuorum Revision}
A Git commit is structurally valid if its object is valid and its hash matches. Trust is usually external policy. A KeyQuorum revision has an explicit trust state machine:
\begin{lstlisting}
CREATED
   v
CHECKED_IN
   v
SIGNED
   v
COUNTERSIGNED / BRIDGE-APPROVED if required
   v
POLICY_EVALUATED
   +--> TRUSTED
   `--> UNTRUSTED / PENDING / DENIED
\end{lstlisting}
A KeyQuorum revision is therefore comparable to a Git commit plus an enforced protected-branch/review/signature policy, except the policy is embedded in the file-security architecture and can depend on HCP topology and physical custody.

\subsection{Author Versus Approver}
Git distinguishes author and committer. KeyQuorum needs an even more explicit split:
\begin{itemize}[leftmargin=1.5em]
  \item \textbf{Revision author}: HCP that created/checked in the content.
  \item \textbf{Content signer}: HCP that attested to the revision.
  \item \textbf{Countersigner}: supervisor/ancestor approving a subordinate action.
  \item \textbf{Bridge approver}: cross-branch authorization mechanism where applicable.
\end{itemize}
These roles can be cryptographically bound and checked against HCP topology.

\subsection{Git Signatures Versus KeyQuorum Signatures}
Git supports signed commits/tags, but signatures are commonly optional and do not encode organizational hierarchy. In KeyQuorum, a cryptographically valid signature can still be insufficient policy evidence.
\begin{lstlisting}
Cryptographically valid signature
        !=
Policy-authorized trusted revision
\end{lstlisting}
For example, \texttt{M.A.1} may sign correctly but still require \texttt{M.A}'s countersignature.

\subsection{Branches and Security-Significant Forks}
Git branches are normal collaboration. KeyQuorum divergence may mean concurrent edits, stale lineage, competing department changes, or an attempted alternate lineage. Both branches should be preserved until policy resolves the conflict.

\subsection{Fast-Forward Semantics}
Git fast-forwards when the current commit is an ancestor of the incoming commit. KeyQuorum can use the same structural test, then add security checks:
\begin{itemize}[leftmargin=1.5em]
  \item existing revision IDs/commitments match,
  \item incoming content proof is valid,
  \item HCP relationship was valid under the bound topology generation,
  \item required countersignatures/bridge/quorum proofs are satisfied,
  \item history extension validates, and
  \item local file policy permits advancement.
\end{itemize}

\subsection{Merge Semantics}
KeyQuorum should allow explicit multi-parent revisions for genuine concurrent edits:
\begin{lstlisting}
R2A(*@\revmark@*) ----+
        +--> R3(*@\revmark@*)
R2B(*@\revmark@*) ----+
\end{lstlisting}
The merge itself can require approval:
\begin{lstlisting}
MERGE_PROPOSED
parents: [R2A(*@\revmark@*), R2B(*@\revmark@*)]
author: M.A.1

REVISION_SIGNED
actor: M.A.1
revision: R3(*@\revmark@*)

COUNTERSIGNATURE_ADDED
actor: M.A
revision: R3(*@\revmark@*)

POLICY_DECISION
result: TRUSTED
\end{lstlisting}

\subsection{History Rewriting}
Git permits rebase, reset, filter operations, and force-push. That flexibility conflicts with KeyQuorum's chain-of-custody goal. Trusted KeyQuorum history should therefore be append-only. Corrections should be represented as new revocation/correction events and revisions instead of rewriting prior trusted history.

\subsection{Content Addressing and Stable File Identity}
Git object identity changes whenever the committed object changes. KeyQuorum similarly makes each \texttt{revision\_id} a cryptographic hash of the canonical revision object, including its parent revision ID or IDs. KeyQuorum nevertheless keeps the long-lived document identity separate from revision identity:

\begin{lstlisting}
file_id            = stable random document identity
revision_id        = SHA-256(canonical revision object)
parent IDs         = cryptographic links to ancestry
content commitment = protected integrity proof for payload
generated_label    = mandatory human-readable revision label
user_label         = optional commit-message-like description
\end{lstlisting}

This allows one logical file to survive edits, moves, renames, devices, and recipients while every immutable revision remains cryptographically linked to its ancestry.

\subsection{Content Hash Privacy}
Git publicly exposes content-addressed IDs. KeyQuorum already avoids unnecessary unkeyed plaintext hashes for some protected files because a plaintext fingerprint can leak information independently of encryption. Therefore KQTF should not blindly copy Git's globally visible plaintext-hash model. Content commitments may stay sealed within KQTF or use a protected commitment construction appropriate to the threat model.

\subsection{Git Remotes Versus KeyQuorum Sharing}
Git remotes synchronize object graphs; they do not inherently provide recipient confidentiality. KeyQuorum sharing is recipient-bound:
\begin{lstlisting}
select trusted revision
        v
bind file_id + revision_id + history_root
        v
sign transport preimage
        v
seal to recipient X25519 key
        v
KQPB relay / transport
        v
recipient verifies and acknowledges
\end{lstlisting}
This is materially different from fetch/push.

\subsection{Branch Protection Versus KeyQuorum Policy}
Protected branches and required reviews are the closest Git-hosting analogy to KeyQuorum revision trust. KeyQuorum's policy is richer because it may depend on HCP ancestry, direct-parent supervisor relationships, topology generation, bridges, logical quorum, minimum physical devices, password/PIN gates, expiry, and document-specific rules.

\subsection{Reflog Versus Embedded Provenance}
Git reflog is local recovery metadata and can expire or be pruned. KeyQuorum history is intended to be part of the tracked artifact, hash-chained, optionally signed, portable, and security-relevant. It is closer to a chain-of-custody record than a reflog.

\subsection{Tags Versus Trusted Checkpoints}
A signed Git tag can mark a release. KeyQuorum can similarly expose named trusted checkpoints, but acceptance can be policy-derived:
\begin{lstlisting}
R7(*@\revmark@*)
name: Approved FY26 Sales Report
status: TRUSTED
scope: M.A
signed_by: M.A
countersigned_by: M
\end{lstlisting}

\subsection{Capabilities Git Does Not Natively Model}
Git does not natively model destructive file expiry, physical-device quorum, HCP topology, supervisor countersignatures, recipient-sealed delivery, password/PIN outcome history, or the distinction between a working revision and a policy-trusted revision.

\subsection{What KeyQuorum Should Borrow From Git}
\begin{itemize}[leftmargin=1.5em]
  \item immutable revision objects,
  \item explicit parent revision IDs,
  \item a revision DAG rather than forced linear history,
  \item fast-forward detection,
  \item multi-parent merge revisions,
  \item deterministic canonical serialization before hashing/signing,
  \item separation of stable object identity from mutable display names,
  \item rebuildable local indexes, and
  \item graph/diff visualization.
\end{itemize}

\subsection{What KeyQuorum Should Not Borrow From Git}
\begin{itemize}[leftmargin=1.5em]
  \item unrestricted history rewriting,
  \item globally exposed plaintext content hashes as the universal identity primitive,
  \item free-form author/committer text as identity proof,
  \item signatures as optional decoration,
  \item unsealed generic object synchronization,
  \item force-update semantics that erase provenance, or
  \item timestamp-only conflict resolution.
\end{itemize}

\subsection{Concise Relationship to Git}
\begin{quote}
KeyQuorum File History borrows Git's immutable revision graph, parent relationships, fast-forward detection, and explicit merge semantics, but replaces Git's developer-centric trust model with HCP topology, supervisor countersigning, quorum, hardware custody, recipient-sealed transport, expiry, and policy-controlled revision trust.
\end{quote}

\begin{quote}
\textbf{Git answers: ``How did these bytes evolve?''}\\[0.4em]
\textbf{KeyQuorum additionally answers: ``Who was authorized to make that evolution trusted, under which topology, custody, and security policy?''}
\end{quote}

\section{Recommended Data-Structure Refinement}
Because revision lineage and security-event ordering differ, model both explicitly.

\subsection{Revision DAG}
The revision ID is a 32-byte SHA-256 digest of the canonical revision object. Parent revision IDs are therefore also 32-byte hashes.

\begin{lstlisting}
pub struct FileRevision {
    // Computed from canonical serialization; not an independent random ID.
    pub revision_id: [u8; 32],

    // Stable document identity shared by every revision of this file.
    pub file_id: [u8; 16],

    // Empty for genesis, one for a normal edit, multiple for a merge.
    pub parent_revision_ids: Vec<[u8; 32]>,

    pub content_commitment: [u8; 32],

    // Always generated by KeyQuorum:
    // R<normalized-file-name>-<UTC timestamp>-<HCP label>
    pub generated_label: String,

    // Optional human description, analogous to a commit message.
    pub user_label: Option<String>,

    pub author_identity: Option<[u8; 16]>,
    pub author_hcp_label: String,
    pub created_at_utc: String,
    pub topology_generation: u64,
    pub policy_hash: [u8; 32],
}
\end{lstlisting}

Canonical encoding must be deterministic. Parent IDs must be encoded in a deterministic order; for a merge, the caller should preserve semantically meaningful parent ordering and the encoder must serialize that exact order. Length prefixes and domain separation must prevent field-boundary ambiguity.

A normal revision therefore points backward cryptographically:

\begin{lstlisting}
revision bc2810f...
parent   a41c9e7...

revision f93ad44...
parent   bc2810f...
\end{lstlisting}

A merge revision carries multiple parent hashes:

\begin{lstlisting}
revision 9c44d72...
parents:
  - 0f6a3c1...
  - 7be219d...
\end{lstlisting}

\subsection{Security Event Chain}
\begin{lstlisting}
pub struct HistoryEvent {
    pub event_id: [u8; 16],
    pub sequence: u64,
    pub file_id: [u8; 16],
    pub revision_id: Option<[u8; 32]>,
    pub previous_event_hash: [u8; 32],
    pub event_hash: [u8; 32],
    pub occurred_at: String,
    pub event_type: HistoryEventType,
    pub outcome: HistoryOutcome,
    pub details: EventDetails,
}
\end{lstlisting}

The KQTF header should carry both:
\begin{lstlisting}
current_revision_id
latest_trusted_revision_id
revision_graph_root / index
history_root
\end{lstlisting}

\section{Suggested CLI Surface}

Human-facing commands should stay concise. KeyQuorum derives pending revision, parent, conflict, reviewer, and policy state whenever possible.

\begin{lstlisting}
keyquorum file track <path>
keyquorum file status <file>
keyquorum file history <file>
keyquorum file graph <file>
keyquorum file diff <file>
keyquorum file checkout <file>
keyquorum file checkin <file>
keyquorum file merge <file>
keyquorum file review <file>
keyquorum file sign <file>
keyquorum file countersign <file>
keyquorum file finalize <file>
keyquorum file verify <file>
keyquorum file history export <file>
keyquorum file history verify <file>
\end{lstlisting}

For interactive terminals, \texttt{merge} should automatically open the human review UI whenever review is required or a clean auto-merge is awaiting acceptance. \texttt{review} reopens the outstanding review. The CLI should infer the relevant pending revision and conflict state instead of requiring the user to copy long revision IDs into the normal workflow.

Lower-level options such as explicit revision IDs and parent hashes may remain available for scripts, tests, recovery, and advanced diagnostics, but they are not required for ordinary human review.


\section{Revision Label Presentation}

Every revision has three complementary forms of identification:

\begin{enumerate}[leftmargin=1.8em]
  \item \textbf{Authoritative revision ID}: the full cryptographic SHA-256 revision hash.
  \item \textbf{Generated label}: always present, using \texttt{R<normalized-file-name>-<UTC timestamp>-<HCP label>}.
  \item \textbf{Optional user label}: a human-entered description analogous to a Git commit message.
\end{enumerate}

A UI should normally present the optional user label first when present, but must retain the generated label and abbreviated revision hash nearby.

\begin{lstlisting}
Updated sales totals after accounting review
Rsales-report-20261002T143205.482Z-M.A
bc2810f...
\end{lstlisting}

When no optional user label exists:

\begin{lstlisting}
Rsales-report-20261002T143205.482Z-M.A
bc2810f...
\end{lstlisting}

The generated label is not required to be globally unique. The revision hash is authoritative and supplies collision-resistant identity.

\section{Browser Lab Representation}

The browser lab should surface tracked-file history inside the existing \textbf{Activity} page rather than creating a separate History tab. This matches the current lab architecture: \texttt{App.tsx} already exposes an Activity tab, \texttt{ActivityPanel.tsx} renders the current action trace plus the full activity log, and \texttt{src/lab/view.rs} / \texttt{src/lab/state.rs} already provide the Rust-to-WASM activity projection.

The Activity page should be upgraded to support both general lab activity and file-scoped provenance.

Recommended controls are:

\begin{itemize}[leftmargin=1.5em]
  \item \textbf{All Activity}: current behavior.
  \item \textbf{Tracked File}: filter by file ID/name.
  \item \textbf{Revision}: filter or expand by revision ID/generated label.
  \item \textbf{Security}: signatures, countersignatures, quorum, PIN/password, expiry, access denial.
  \item \textbf{Sharing}: share attempts, delivery, acknowledgement, fallback-to-trusted-version.
  \item \textbf{Conflict}: fork detection, auto-merge attempts, human escalation, merge resolution.
\end{itemize}

A tracked-file activity entry may render:

\begin{lstlisting}
AUTO_MERGE_CLEAN
decision.txt

revision:
  Rdecision-20261002T143205.482Z-M.A

parents:
  0f6a3c1...
  7be219d...

outcome:
  merged content generated

trust:
  PENDING - signature required
\end{lstlisting}

For a conflict requiring a human:

\begin{lstlisting}
CONTENT_CONFLICT_DETECTED
decision.txt

M.A.2 and M.S.1 both modified line 66

automatic merge:
  REQUIRES_HUMAN

review assigned:
  M.A

selection:
  PRIOR_NEUTRAL_OWNER
\end{lstlisting}

The page should also offer an expandable revision graph for the selected file. A fork must be visually obvious and show which heads are trusted, pending, denied, or unresolved.

The lab should demonstrate at least three history scenarios:

\begin{enumerate}[leftmargin=1.8em]
  \item \textbf{Unsigned fallback}: a newer unsigned revision exists, but a request receives the prior trusted revision.
  \item \textbf{Automatic clean merge}: two owners edit different lines; KeyQuorum produces a multi-parent merge revision automatically, then shows its pending/trusted policy state.
  \item \textbf{Human conflict escalation}: two owners modify the same line differently; the Activity page shows automatic merge failure, prior-owner/seniority resolution selection, and the final merge decision.
\end{enumerate}

No React component should independently recalculate ancestry, merge eligibility, signature validity, or trust. The browser lab continues to invoke the real Rust/CLI behavior compiled to WASM and renders the returned activity/history state.

The browser lab may visualize the same changed-range provenance inside Activity, but the native CLI review behavior remains part of the Rust CLI/TUI path rather than React.


\section{Suggested Module Changes}

{\small
\setlength{\tabcolsep}{5pt}
\renewcommand{\arraystretch}{1.12}
\begin{tabularx}{\textwidth}{
  >{\raggedright\arraybackslash}p{0.33\textwidth}
  >{\raggedright\arraybackslash}X
}
\toprule
\textbf{Module} & \textbf{History Integration / Upgrade} \\
\midrule
\texttt{src/file\_history.rs} & New orchestration surface for KQTF, revision DAG, provenance, and history APIs; delegates existing security decisions to existing modules. \\
\texttt{src/file\_history/container.rs} & New KQTF encoding/decoding and container integrity. \\
\texttt{src/file\_history/event.rs} & New event types, canonical serialization, and hash-chain logic. \\
\texttt{src/file\_history/revision.rs} & New revision DAG, parent hashes, commitments, and multi-parent revision construction. \\
\texttt{src/file\_history/policy.rs} & Compose existing authority/signature/bridge/quorum results into revision trust, finalization state, and delivery selection; do not reproduce those checks. \\
\texttt{src/file\_history/topology.rs} & Thin adapter only if needed; shared ancestry algorithms should be promoted from existing topology helpers rather than duplicated here. \\
\texttt{src/file\_history/merge.rs} & New deterministic three-way content merge, fork/conflict classification, auto-merge outcome, resolver selection, and escalation; calls shared topology/trust modules. \\
\texttt{src/file\_history/verify.rs} & Verify KQTF/history/revision graph and call existing signing/bridge verifiers for cryptographic proofs. \\
\texttt{src/authority.rs} & Upgrade to expose shared HCP relationship, ancestry distance, direct-parent, and lowest-common-ancestor primitives used by existing authority flows and File History. \\
\texttt{src/private\_bridge.rs} & Reuse current bridge roster, supervisor metadata, privacy, and bridge signature semantics; migrate/re-export topology helpers if centralized. \\
\texttt{src/signing.rs} & Add revision/countersign/history domain preimages; retain existing \texttt{sign}, \texttt{verify\_signature}, and bridge-signature implementations. \\
\texttt{src/file\_delivery.rs} & Version/extend existing delivery and acknowledgement structures to carry tracked-file metadata; do not add a parallel delivery protocol. \\
\texttt{src/envelope.rs} & Add history message kinds and reuse existing KQPB/KQXB framing and codec helpers. \\
\texttt{src/quorum.rs} & Emit tracked history around the existing reconstruction/device/parent-approval/decrypt path; do not duplicate quorum evaluation. \\
\texttt{src/key\_tree.rs} & Reuse existing recursive tree/reconstruction data for quorum and topology evidence where applicable; no history-specific secret-sharing implementation. \\
\texttt{src/device.rs} & Reuse existing custody policy and presented-device accounting for history outcomes. \\
\texttt{src/locked\_files.rs} & Reuse existing password/expiry implementation and add tracked-history hooks/tombstone behavior. \\
\texttt{src/pin.rs} & Keep PIN verification generic; callers emit history outcomes after the existing verifier returns. \\
\texttt{src/sharing.rs} & Instrument existing create/redeem/revoke paths with history events; preserve current token-hash, expiry, and concurrency logic. \\
\texttt{src/transfer.rs} & Reuse stable cryptographic identity and possession semantics for actor attribution; do not conflate key provenance with file ownership history. \\
\texttt{src/export.rs} & Upgrade existing KQXB export to preserve KQTF/history snapshots rather than introducing a separate export channel. \\
\texttt{src/db/schema.sql} & Add rebuildable tracked-file/revision/history indexes; authoritative history remains in KQTF. \\
\texttt{src/db/mod.rs} & Extend the existing migration path for new cache/index tables; no second migration framework. \\
\texttt{src/lab/view.rs} & Extend the existing \texttt{ActivityView} with optional file/revision/history/conflict fields. \\
\texttt{src/lab/state.rs} & Project real CLI/history results into the existing activity log; do not reimplement security or merge policy in lab state. \\
\texttt{lab/src/api/types.ts} & Mirror the upgraded Rust \texttt{ActivityView}; Rust remains source of truth. \\
\texttt{lab/src/components/ActivityPanel.tsx} & Upgrade the existing Activity page with file filters, history entries, conflict/merge status, and revision-graph expansion; do not add a separate History panel. \\
\texttt{src/cli.rs / CLI command layer} & Upgrade merge/review commands so interactive terminals automatically enter review mode; reuse the same File History merge/revision APIs used by non-interactive commands. \\
\texttt{src/file\_history/merge.rs} & Expose changed-range provenance and merge-review view models; do not move author inference or merge policy into the terminal renderer. \\
\texttt{native TUI renderer} & Add ratatui/crossterm presentation for color-coded changed lines, mouse/keyboard provenance inspection, review actions, signing, and finalization. No security logic lives here. \\
\bottomrule
\end{tabularx}
}


\section{Implementation Phases}
\begin{enumerate}[leftmargin=1.8em]
  \item \textbf{Shared-helper refactor}: promote/reuse ancestry/topology primitives and establish the ``upgrade, do not duplicate'' module boundaries before File History depends on them.
  \item \textbf{Container foundation}: KQTF framing, stable file IDs, canonical serialization, event hashing.
  \item \textbf{Revision DAG}: parent hashes, generated/user labels, current/trusted heads, fast-forward/fork detection.
  \item \textbf{Content signatures}: add revision/history domain preimages to the existing signing module; keep transport signatures separate.
  \item \textbf{HCP policy}: compose shared ancestry, scope, topology generation, supervisor countersignature, and bridge state.
  \item \textbf{Automatic conflict resolution}: deterministic three-way text merge, clean-merge proposals, policy gating, and explicit fallback to human resolver selection.
  \item \textbf{Human conflict escalation}: prior-neutral-owner selection, seniority, sector prominence, common-ancestor escalation, and bridge-assisted routing.
  \item \textbf{Delivery integration}: version/extend existing file delivery/ack with tracked IDs/history roots and latest-trusted fallback.
  \item \textbf{Password/PIN/quorum/expiry}: emit safe history outcomes around existing modules and implement tracked tombstones.
  \item \textbf{History export}: HistorySnapshot through the existing KeyQuorum envelope/export family.
  \item \textbf{Database indexes/migration}: add rebuildable indexes using the current DB migration system.
  \item \textbf{Lab Activity integration}: extend the existing Rust ActivityView and ActivityPanel with tracked-file filters, revision graphs, auto-merge, and conflict-resolution events.
\end{enumerate}


\section{Testing Requirements}

Automated tests should cover tracking activation; stable identity across rename/move; deterministic revision IDs; parent-hash linkage; single- and multi-parent revisions; current versus trusted selection; supervisor countersigning; scope-owner self-sign; topology-generation binding; cross-branch and bridge rules; private bridge nondisclosure; content versus transport signatures; unsigned latest-revision fallback; password/PIN logging without leakage; quorum/physical-device policy; expiry tombstones; post-expiry access logging; tamper detection; fast-forward; equivalent-head collapse; one-sided three-way merge; non-overlapping automatic text merge; identical overlapping edits; overlapping conflicting edits; unsupported/binary content fallback; policy conflict; auto-merge not inheriting signatures; merge-review state; finalization signatures; latest-final-head tracking; colored diff classification for added/removed/unchanged lines; changed-lines-first ordering; hover/keyboard provenance lookup; multi-contributor line provenance; concise merge/review command inference; TUI/non-interactive behavior parity; prior-neutral-owner review selection; HCP seniority fallback; sector-prominence tie-breaking; least-common-ancestor escalation; bridge-assisted escalation; unresolved-conflict handling; explicit merge revisions; history export; cache rebuild; Activity-page projection; and no fabricated attribution.

Regression tests must additionally prove that File History uses the same existing security primitives rather than parallel implementations. Examples include confirming that revision signature verification calls the same strict Ed25519 verifier, tracked unlocks use the same quorum/device policy path, PIN behavior matches the existing PIN verifier, share redemption preserves current atomic use-count behavior, and lab actions still execute through the same real CLI/WASM path.


\section{Architectural Summary}
KeyQuorum File History is not merely an audit log and not merely a version-control system. It is a portable, cryptographically verifiable chain-of-custody and revision-trust system built around:
\begin{center}
\textbf{Stable File Identity + Revision DAG + HCP Topology + Supervisor Countersigning + Bridge Authorization + Quorum Decisions + Content Signatures + Transport Signatures + Append-Only Security History + Policy-Controlled Revision Selection}
\end{center}

Git provides the structural analogy for immutable revisions, parent relationships, branches, fast-forwards, and merges. KeyQuorum extends that structure into a security architecture where a revision does not become trusted merely because it exists; it becomes trusted when the applicable cryptographic, organizational, custody, and file policy is satisfied.

1 & Covered & KQTF answers the provenance, trust, sharing, gate, expiry, fork and conflict questions through its revision graph and event chain. Cross-branch decisions name the privacy-safe approval alternative without identifying a private bridge. \\
